\begingroup\setbox0=\vbox{\section*{Direct and Inverse Proportion}}\endgroup
\chapterhead{7}{Direct and Inverse Proportion}

\shead{Two proportional quantities}
\begin{knowbox}
\textbf{Direct proportion.} Two quantities are directly proportional if, when one is multiplied (or divided) by some number, the other is multiplied (or divided) by the same number. \emph{Example:} metres of cloth and its cost (at a fixed price).\par
\textbf{Inverse proportion.} Two quantities are inversely proportional if, when one is multiplied by some number, the other is divided by that number; their \emph{product} stays the same. \emph{Example:} the number of workers and the number of days needed for the same job.\par
\smallskip
Two common methods:
\begin{itemize}[leftmargin=1.4em,itemsep=3pt]
\item \textbf{The unitary method:} find the value for $1$ unit (1 metre, 1 worker, 1 day,\ldots), then multiply or divide.
\item \textbf{The ratio method:} find how many times the first quantity grows (or shrinks), then change the second one accordingly.
\end{itemize}
For ``workers -- days'' problems, think of the job as \emph{(workers) $\times$ (days)} ``worker-days''; the number of worker-days needed does not change.
\end{knowbox}
\gap
\begin{exbox}{\ExWord\ 1}
$5$ m of cloth cost $450\,000$ dong. How much do $15$ m of the same cloth cost?
\guide
\textbf{Method 1 (unitary).} $1$ m costs $450\,000:5=90\,000$ dong, so $15$ m cost $90\,000\times15=1\,350\,000$ dong.\par
\textbf{Method 2 (ratio).} $15$ m is $15:5=3$ times $5$ m, so the cost is $3$ times as much: $450\,000\times3=1\,350\,000$ dong.
\dapso{$1\,350\,000$ dong.}
\end{exbox}
\gap
\begin{exbox}{\ExWord\ 2}
A team of $12$ workers builds a stretch of road in $10$ days. How many workers are needed to finish it in $8$ days? (Everyone works at the same rate.)
\guide
Workers and days are inversely proportional.\par One worker alone would need $10\times12=120$ days, so the job takes $120$ worker-days.\par
To finish in $8$ days we need $120:8=15$ workers.
\dapso{$15$ workers.}
\end{exbox}
\gap
\begin{exbox}{\ExWord\ 3}
A kitchen has enough rice for $120$ people for $20$ days. After $5$ days, $30$ more people arrive. For how many more days will the remaining rice last? (Everyone eats the same amount.)
\guide
After $5$ days, the rice left would feed $120$ people for $20-5=15$ days.\par
That rice would feed $1$ person for $15\times120=1\,800$ days.\par
There are now $120+30=150$ people, so the rice lasts $1\,800:150=12$ days.
\dapso{$12$ days.}
\end{exbox}
\gap
\begin{exbox}{\ExWord\ 4}
$10$ workers repair $400$ m of road in $5$ days. How many metres can $12$ workers repair in $6$ days? (Everyone works at the same rate.)
\guide
There are three quantities; reduce to ``$1$ worker in $1$ day''.\par
In $5$ days $1$ worker repairs $400:10=40$ m; in $1$ day, $40:5=8$ m.\par
$12$ workers in $6$ days repair $8\times12\times6=576$ m.
\dapso{$576$ m.}
\end{exbox}
\dividerbar{Practice}
\begin{probbox}{1}{\relax}
$6$ identical boxes hold $72$ pens altogether. How many pens do $15$ such boxes hold?
\end{probbox}
\probsol{
Each box holds $72:6=12$ pens, so $15$ boxes hold $12\times15=180$ pens.
\dapso{$180$ pens.}}
\gap
\begin{probbox}{2}{\relax}
A car uses $9$ litres of petrol for $120$ km. How much petrol does it use for $200$ km?
\end{probbox}
\probsol{
$40$ km uses $9:3=3$ litres (because $120=40\times3$).\par $200$ km is $200:40=5$ times $40$ km, so it uses $3\times5=15$ litres.
\dapso{$15$ litres.}}
\gap
\begin{probbox}{3}{\relax}
$15$ outfits use $45$ m of cloth. How many such outfits can be made from $72$ m?
\end{probbox}
\probsol{
Each outfit needs $45:15=3$ m, so $72$ m make $72:3=24$ outfits.
\dapso{$24$ outfits.}}
\gap
\begin{probbox}{4}{\relax}
With enough money for $24$ notebooks at $6\,000$ dong each, how many notebooks at $8\,000$ dong each can you buy?
\end{probbox}
\probsol{
The money is $6\,000\times24=144\,000$ dong, which buys $144\,000:8\,000=18$ notebooks.
\dapso{$18$ notebooks.}}
\gap
\begin{probbox}{5}{\relax}
$15$ workers finish a job in $12$ days. How long would $20$ workers take? (Everyone works at the same rate.)
\end{probbox}
\probsol{
The job takes $15\times12=180$ worker-days, so $20$ workers need $180:20=9$ days.
\dapso{$9$ days.}}
\gap
\begin{probbox}{6}{\relax}
A kitchen has rice for $80$ people for $15$ days. How long would it last $100$ people?
\end{probbox}
\probsol{
The rice would feed $1$ person for $15\times80=1\,200$ days, so $100$ people for $1\,200:100=12$ days.
\dapso{$12$ days.}}
\gap
\begin{probbox}{7}{\relax}
A summer camp of $60$ pupils has food for $30$ days. After $10$ days, $20$ pupils go home early. For how many more days will the food last those who stay?
\end{probbox}
\probsol{
The food left would feed $60$ pupils for $20$ days, i.e.\par $60\times20=1\,200$ pupil-days.\par Now $40$ pupils remain, so it lasts $1\,200:40=30$ days.
\dapso{$30$ days.}}
\gap
\begin{probbox}{8}{\relax}
$8$ people plant $96$ trees in $3$ hours. How many trees do $10$ people plant in $4$ hours? (Everyone works at the same rate.)
\end{probbox}
\probsol{
$1$ person plants $96:8:3=4$ trees per hour.\par $10$ people in $4$ hours: $4\times10\times4=160$ trees.
\dapso{$160$ trees.}}
\gap
\begin{probbox}{9}{\relax}
$3$ workers finish a job in $8$ days. How many more workers are needed to finish it $2$ days sooner? (Everyone works at the same rate.)
\end{probbox}
\probsol{
The job is $3\times8=24$ worker-days.\par In $8-2=6$ days we need $24:6=4$ workers, i.e.\par $4-3=1$ more.
\dapso{$1$ worker.}}
\gap
\begin{probbox}{10}{\relax}
Every $100$ kg of paddy gives $62$ kg of rice. How many kilograms of rice does $2\,500$ kg of paddy give?
\end{probbox}
\probsol{
$2\,500$ kg is $25$ times $100$ kg, so the rice is $62\times25=1\,550$ kg.
\dapso{$1\,550$ kg.}}
\gap
\begin{probbox}{11}{\relax}
$5$ identical pumps move $360$ m$^3$ of water in $4$ hours. How much do $8$ such pumps move in $3$ hours?
\end{probbox}
\probsol{
$1$ pump in $1$ hour: $360:5:4=18$ m$^3$.\par $8$ pumps in $3$ hours: $18\times8\times3=432$ m$^3$.
\dapso{$432$ m$^3$.}}
\gap
\begin{probbox}{12}{\relax}
A workshop of $45$ workers plans to finish an order in $20$ days. After $4$ days, $15$ more workers join. How many days early is the order finished?
\end{probbox}
\probsol{
The rest of the work is $45$ workers for $16$ days, i.e.\par $45\times16=720$ worker-days.\par With $60$ workers it takes $720:60=12$ days.\par
In total $4+12=16$ days, which is $20-16=4$ days early.
\dapso{$4$ days.}}
\gap
\begin{probbox}{13}{\relax}
A kitchen has rice for $60$ people for $24$ days. After $4$ days more people arrive, and the rice left lasts only $16$ more days. How many people arrived?
\end{probbox}
\probsol{
The rice left is $60\times20=1\,200$ person-days.\par Over $16$ days that feeds $1\,200:16=75$ people.\par
So $75-60=15$ people arrived.
\dapso{$15$ people.}}
\gap
\begin{probbox}{14}{\relax}
$4$ kg of apples cost $120\,000$ dong. How much do $7$ kg cost?
\end{probbox}
\probsol{
$1$ kg: $120\,000:4=30\,000$ dong.\par $7$ kg: $30\,000\times7=210\,000$ dong.
\dapso{$210\,000$ dong.}}
\gap
\begin{probbox}{15}{\relax}
$12$ notebooks cost $96\,000$ dong. How much do $20$ cost?
\end{probbox}
\probsol{
One: $96\,000:12=8\,000$ dong.\par $20$: $8\,000\times20=160\,000$ dong.
\dapso{$160\,000$ dong.}}
\gap
\begin{probbox}{16}{\relax}
A car covers $135$ km in $3$ hours. How far does it go in $5$ hours? (Same distance every hour.)
\end{probbox}
\probsol{
One hour: $135:3=45$ km.\par $5$ hours: $45\times5=225$ km.
\dapso{$225$ km.}}
\gap
\begin{probbox}{17}{\relax}
$8$ workers finish a job in $6$ days. How long would $12$ workers take? (Same rate.)
\end{probbox}
\probsol{
The job is $8\times6=48$ worker-days.\par $12$ workers need $48:12=4$ days.
\dapso{$4$ days.}}
\gap
\begin{probbox}{18}{\relax}
Reading $15$ pages a day, Ha finishes a book in $12$ days. How long would she take at $20$ pages a day?
\end{probbox}
\probsol{
The book has $15\times12=180$ pages.\par Days: $180:20=9$.
\dapso{$9$ days.}}
\gap
\begin{probbox}{19}{\relax}
A kitchen has rice for $50$ people for $18$ days. How long would it last $60$ people?
\end{probbox}
\probsol{
The rice feeds $1$ person for $18\times50=900$ days.\par $60$ people: $900:60=15$ days.
\dapso{$15$ days.}}
\gap
\begin{probbox}{20}{\relax}
$6$ identical sewing machines make $600$ shirts in $5$ days. How many do $10$ machines make in $3$ days?
\end{probbox}
\probsol{
$1$ machine in $1$ day: $600:6:5=20$ shirts.\par $10$ machines in $3$ days: $20\times10\times3=600$ shirts.
\dapso{$600$ shirts.}}
\gap
\begin{probbox}{21}{\relax}
$6$ m of cable cost $54\,000$ dong. How many metres can $360\,000$ dong buy?
\end{probbox}
\probsol{
$1$ m costs $54\,000:6=9\,000$ dong.\par $360\,000:9\,000=40$ m.
\dapso{$40$ m.}}
\gap
\begin{probbox}{22}{\relax}
A tap delivers $60$ litres in $5$ minutes. How much does it deliver in $1$ hour?
\end{probbox}
\probsol{
$1$ hour $=60$ minutes, which is $60:5=12$ times $5$ minutes.\par Water: $60\times12=720$ litres.
\dapso{$720$ litres.}}
\gap
\begin{probbox}{23}{\relax}
$6$ identical lorries carry $72$ tonnes. How much do $15$ such lorries carry?
\end{probbox}
\probsol{
Each lorry: $72:6=12$ t.\par $15$ lorries: $12\times15=180$ t.
\dapso{$180$ t.}}
\gap
\begin{probbox}{24}{\relax}
A team of $20$ plans to finish a job in $15$ days. How many more workers are needed to finish in $12$ days? (Same rate.)
\end{probbox}
\probsol{
The job is $20\times15=300$ worker-days.\par In $12$ days we need $300:12=25$ workers.\par So $25-20=5$ more.
\dapso{$5$ workers.}}
\gap
\begin{probbox}{25}{\relax}
$9$ people plan to finish a job in $12$ days. After $4$ days, $3$ of them stop. How many days late is the job finished?
\end{probbox}
\probsol{
Work left: $9\times(12-4)=72$ worker-days.\par $6$ people need $72:6=12$ days.\par Total: $4+12=16$ days, which is $16-12=4$ days late.
\dapso{$4$ days.}}
\gap
\begin{probbox}{26}{\relax}
Every $30$ kg of paddy gives $18$ kg of rice. How much rice does $1$ tonne of paddy give?
\end{probbox}
\probsol{
$1$ t $=1\,000$ kg.\par $1$ kg of paddy gives $18:30=0.6$ kg of rice.\par $1\,000$ kg give $0.6\times1\,000=600$ kg.
\dapso{$600$ kg.}}
\gap
\begin{probbox}{27}{\relax}
$5$ identical sacks of rice weigh $250$ kg. How much do $12$ sacks weigh?
\end{probbox}
\probsol{
One sack: $250:5=50$ kg.\par $12$ sacks: $50\times12=600$ kg.
\dapso{$600$ kg.}}
\gap
\begin{probbox}{28}{\relax}
A school pays $9\,000\,000$ dong for $15$ desk sets. How much more money is needed to buy $25$ sets?
\end{probbox}
\probsol{
One set: $9\,000\,000:15=600\,000$ dong.\par $25$ sets: $600\,000\times25=15\,000\,000$ dong.\par Extra: $15\,000\,000-9\,000\,000=6\,000\,000$ dong.
\dapso{$6\,000\,000$ dong.}}
\gap
\begin{probbox}{29}{\relax}
Carrying $4$ tonnes a trip, a lorry needs $15$ trips. How many trips are needed at $6$ tonnes a trip?
\end{probbox}
\probsol{
Load: $4\times15=60$ t.\par Trips: $60:6=10$.
\dapso{$10$ trips.}}
\gap
\begin{probbox}{30}{\relax}
$14$ people eat $21$ kg of rice in $3$ days. How much do $20$ people eat in $6$ days? (Same amount each.)
\end{probbox}
\probsol{
$1$ person for $1$ day: $21:14:3=0.5$ kg.\par $20$ people for $6$ days: $0.5\times20\times6=60$ kg.
\dapso{$60$ kg.}}
\dividerbar{Challenge -- Extension}
\begin{probbox}{31}{\stars{1}}
A unit has rice for $30$ days. After $10$ days, $40$ more people join, and the rice left lasts only $15$ more days. How many people were in the unit at first?
\end{probbox}
\probsol{
The rice left feeds the original group for $20$ days, or the original group plus $40$ for $15$ days.\par
Since $20:15=\dfrac43$, the new group is $\dfrac43$ of the original; so $40$ people are $\dfrac13$ of the original group.\par
Original group: $40\times3=120$ people.
\dapso{$120$ people.}}
\gap
\begin{probbox}{32}{\stars{1}}
A team plans to finish a project in $15$ days. After $5$ days, $6$ workers leave, so the team needs $2$ days more than planned. How many workers were there at first?
\end{probbox}
\probsol{
The rest of the work was planned for $10$ days but took $12$.\par The days are $\dfrac{12}{10}=\dfrac65$ times as many, so the workers left are $\dfrac56$ of the original team.\par
So $6$ workers are $1-\dfrac56=\dfrac16$ of the team, which had $6\times6=36$ workers.
\dapso{$36$ workers.}}
\gap
\begin{probbox}{33}{\stars{1}}
$12$ people working $8$ hours a day finish a job in $10$ days. How many days do $16$ people working $6$ hours a day need? (Everyone works at the same rate.)
\end{probbox}
\probsol{
The job is $12\times8\times10=960$ worker-hours.\par Each day $16$ people put in $16\times6=96$ worker-hours.\par
Days needed: $960:96=10$.
\dapso{$10$ days.}}
\gap
\begin{probbox}{34}{\stars{1}}
Two gear wheels mesh together. The big one has $48$ teeth and the small one $18$. When the big wheel turns $15$ times, how many times does the small one turn?
\end{probbox}
\probsol{
The wheels mesh, so the same number of teeth pass the contact point on both: $48\times15=720$.\par
The small wheel turns $720:18=40$ times.
\note{Number of teeth and number of turns are inversely proportional.}
\dapso{$40$ turns.}}
\gap
\begin{probbox}{35}{\stars{2}}
``A hen and a half lays an egg and a half in a day and a half.'' How many eggs do $6$ such hens lay in $6$ days?
\end{probbox}
\probsol{
$1.5$ hens lay $1.5$ eggs in $1.5$ days, so $1.5$ hens lay $1$ egg in $1$ day.\par
$6$ hens are $4$ times $1.5$ hens, so they lay $4$ eggs a day, and $4\times6=24$ eggs in $6$ days.
\dapso{$24$ eggs.}}
\gap
\begin{probbox}{36}{\stars{2}}
With a sum of money, Ms Lan can buy $20$ kg of type I rice or $30$ kg of type II rice. If she spends it on equal amounts of both types, how many kilograms of each does she get?
\end{probbox}
\probsol{
$1$ kg of type I costs $\dfrac1{20}$ of the money, $1$ kg of type II costs $\dfrac1{30}$.\par One kilogram of each costs $\dfrac1{20}+\dfrac1{30}=\dfrac1{12}$ of the money.\par
So she gets $12$ kg of each.
\dapso{Each type: $12$ kg.}}
\gap
\begin{probbox}{37}{\stars{1}}
A team plans to finish a project in $20$ days. With $10$ more workers it would finish $4$ days sooner. How many workers are in the team? (Same rate.)
\end{probbox}
\probsol{
The days drop from $20$ to $16$, i.e.\par $\dfrac{16}{20}=\dfrac45$ of the time, so the team must grow by $\dfrac54$.\par So $10$ workers are $\dfrac14$ of the team.\par Team: $10\times4=40$.
\dapso{$40$ workers.}}
\gap
\begin{probbox}{38}{\stars{1}}
A kitchen has rice for $100$ people for $30$ days. After $6$ days, $20$ people leave. How long will the rest of the rice last?
\end{probbox}
\probsol{
The rice left is $100\times24=2\,400$ person-days.\par $80$ people remain: $2\,400:80=30$ days.
\dapso{$30$ days.}}
\gap
\begin{probbox}{39}{\stars{1}}
A kitchen has rice for $24$ days. After $4$ days, $20$ more people arrive and the rest lasts only $15$ days. How many people were there at first?
\end{probbox}
\probsol{
The rice left feeds the original group for $20$ days, or the group plus $20$ for $15$ days.\par $20:15=\dfrac43$, so the new group is $\dfrac43$ of the old; the $20$ newcomers are $\dfrac13$ of it.\par At first: $60$ people.
\dapso{$60$ people.}}
\gap
\begin{probbox}{40}{\stars{1}}
$6$ identical taps fill a tank in $4$ hours. All $6$ run for $1$ hour, then $2$ are turned off. How much longer do the others need?
\end{probbox}
\probsol{
Filling takes $6\times4=24$ tap-hours.\par The first hour gives $6$; $18$ remain.\par $4$ taps need $18:4=4.5$ hours $=4$ h $30$ min.
\dapso{$4$ h $30$ min.}}
\gap
\begin{probbox}{41}{\stars{1}}
Gears A, B, C mesh in a chain (A with B, B with C). A has $30$ teeth, B $20$ and C $45$. When A turns $12$ times, how many times does C turn?
\end{probbox}
\probsol{
The same number of teeth pass each contact: $30\times12=360$.\par B turns $360:20=18$ times; C also has $360$ teeth pass, so it turns $360:45=8$ times.
\dapso{$8$ turns.}}
\gap
\begin{probbox}{42}{\stars{1}}
A clock gains $3$ minutes a day. After how many days will it be $1$ hour fast?
\end{probbox}
\probsol{
$1$ hour $=60$ minutes.\par Days: $60:3=20$.
\dapso{$20$ days.}}
\gap
\begin{probbox}{43}{\stars{1}}
Walking at $4$ km an hour, Mai takes $30$ minutes to get to school. How long would she take cycling at $12$ km an hour?
\end{probbox}
\probsol{
The distance is fixed; going $12:4=3$ times as fast takes a third of the time.\par Time: $30:3=10$ minutes.
\dapso{$10$ minutes.}}
\gap
\begin{probbox}{44}{\stars{1}}
$4$ people working $8$ hours a day finish a job in $6$ days. How long would $6$ people working $4$ hours a day take?
\end{probbox}
\probsol{
The job is $4\times8\times6=192$ worker-hours.\par Each day $6$ people give $6\times4=24$.\par Days: $192:24=8$.
\dapso{$8$ days.}}
\gap
\begin{probbox}{45}{\stars{1}}
$18$ lorries each carry $45$ sacks. If only $15$ lorries carry the whole load, how many extra sacks must each take?
\end{probbox}
\probsol{
Sacks: $45\times18=810$.\par Each lorry: $810:15=54$; that is $54-45=9$ more.
\dapso{$9$ sacks.}}
\gap
\begin{probbox}{46}{\stars{2}}
``Three cats catch three mice in three minutes.'' How long do $100$ cats take to catch $100$ mice? (All cats are equally quick.)
\end{probbox}
\probsol{
Three cats catch three mice in $3$ minutes, so each cat catches $1$ mouse in $3$ minutes.\par $100$ cats each catching $1$ mouse also take $3$ minutes.
\dapso{$3$ minutes.}}
\gap
\begin{probbox}{47}{\stars{2}}
A sum of money buys $15$ kg of apples or $10$ kg of pears. After buying $6$ kg of apples, how many kilograms of pears can the rest buy?
\end{probbox}
\probsol{
$1$ kg of apples costs $\dfrac1{15}$ of the money; $6$ kg cost $\dfrac6{15}=\dfrac25$.\par The remaining $\dfrac35$ buys $10\times\dfrac35=6$ kg of pears.
\dapso{$6$ kg of pears.}}
\gap
\begin{probbox}{48}{\stars{2}}
$3$ kg of oranges and $2$ kg of apples cost $175\,000$ dong, and apples cost twice as much per kilogram. Find each price.
\end{probbox}
\probsol{
$2$ kg of apples cost the same as $4$ kg of oranges, so the bill is the price of $3+4=7$ kg of oranges.\par Oranges: $175\,000:7=25\,000$ dong/kg; apples: $50\,000$ dong/kg.
\dapso{Oranges $25\,000$, apples $50\,000$ dong/kg.}}
\gap
\begin{probbox}{49}{\stars{2}}
A team of $12$ plans to build $360$ m of dyke in $6$ days. After $2$ days, $4$ more people join. How many days early do they finish?
\end{probbox}
\probsol{
Each person builds $360:12:6=5$ m a day.\par In $2$ days they build $12\times5\times2=120$ m; $240$ m remain.\par $16$ people build $80$ m a day, needing $240:80=3$ days.\par Total $5$ days: $1$ day early.
\dapso{$1$ day.}}
\gap
\begin{probbox}{50}{\stars{2}}
The same money buys $8$ pens or $12$ notebooks. A pen costs $2\,000$ dong more than a notebook. How much money is it?
\end{probbox}
\probsol{
Buying $8$ pens instead of $8$ notebooks costs $8\times2\,000=16\,000$ dong more.\par That extra is the price of $12-8=4$ notebooks, so a notebook is $4\,000$ dong.\par Money: $4\,000\times12=48\,000$ dong.
\dapso{$48\,000$ dong.}}
\clearpage
\titlehead
\begingroup\setbox0=\vbox{\section*{Work Problems}}\endgroup
\chapterhead{8}{Work Problems}

\shead{Treat the whole job as one unit}
\begin{knowbox}
\begin{itemize}[leftmargin=1.4em,itemsep=3pt]
\item Treat \textbf{the whole job} (or \textbf{a full tank}) as $1$ unit.
\item If someone finishes the job in $n$ hours, then \textbf{each hour} they do $\dfrac1n$ of the job.
\item When several people \textbf{work together}, the work done in $1$ hour is the \emph{sum} of their separate parts.
\item An \textbf{outlet pipe} empties the tank, so its part is \emph{subtracted}.
\item Time needed $=$ (work to be done) $:$ (work done in $1$ hour).
\end{itemize}
\end{knowbox}
\gap
\begin{exbox}{\ExWord\ 1}
The first person can finish a job alone in $6$ hours, the second in $3$ hours. How long will it take them working together?
\guide
Each hour the first person does $\dfrac16$ of the job and the second $\dfrac13$.\par
Together they do $\dfrac16+\dfrac13=\dfrac12$ of the job each hour.\par
Time together: $1:\dfrac12=2$ hours.
\dapso{$2$ hours.}
\end{exbox}
\gap
\begin{exbox}{\ExWord\ 2}
Two people together finish a job in $12$ days. The first alone would take $20$ days. How long would the second take alone?
\guide
Each day both together do $\dfrac1{12}$ of the job and the first does $\dfrac1{20}$.\par
So the second does $\dfrac1{12}-\dfrac1{20}=\dfrac5{60}-\dfrac3{60}=\dfrac1{30}$ each day, and would take $30$ days alone.
\dapso{$30$ days.}
\end{exbox}
\gap
\begin{exbox}{\ExWord\ 3}
Pipe I fills a tank in $4$ hours, pipe II in $6$ hours. The tank is empty; both pipes run for $1$ hour, then pipe I is turned off. How much longer must pipe II run to fill the tank?
\guide
In $1$ hour both pipes fill $\dfrac14+\dfrac16=\dfrac5{12}$ of the tank, leaving $1-\dfrac5{12}=\dfrac7{12}$.\par
Pipe II needs $\dfrac7{12}:\dfrac16=\dfrac72$ hours $=3$ hours $30$ minutes.
\dapso{$3$ hours $30$ minutes.}
\end{exbox}
\gap
\begin{exbox}{\ExWord\ 4}
An inlet fills a tank in $3$ hours; an outlet empties a full tank in $5$ hours. Starting empty with both open, how long until the tank is full?
\guide
Each hour the water rises by $\dfrac13-\dfrac15=\dfrac2{15}$ of the tank.\par
Time to fill: $1:\dfrac2{15}=\dfrac{15}2$ hours $=7$ hours $30$ minutes.
\dapso{$7$ hours $30$ minutes.}
\end{exbox}
\dividerbar{Practice}
\begin{probbox}{1}{\relax}
An sweeps the schoolyard in $30$ minutes, Binh in $20$ minutes. How long does it take them together?
\end{probbox}
\probsol{
Each minute together they sweep $\dfrac1{30}+\dfrac1{20}=\dfrac5{60}=\dfrac1{12}$ of the yard, so they take $12$ minutes.
\dapso{$12$ minutes.}}
\gap
\begin{probbox}{2}{\relax}
Pipe I fills a tank in $2$ hours, pipe II in $3$ hours. Starting empty with both open, how long until it is full?
\end{probbox}
\probsol{
Each hour both fill $\dfrac12+\dfrac13=\dfrac56$ of the tank.\par Time: $1:\dfrac56=\dfrac65$ hours $=1$ hour $12$ minutes.
\dapso{$1$ hour $12$ minutes.}}
\gap
\begin{probbox}{3}{\relax}
Two workers together finish a job in $6$ hours. The first alone takes $10$ hours. How long does the second take alone?
\end{probbox}
\probsol{
Each hour the second does $\dfrac16-\dfrac1{10}=\dfrac1{15}$ of the job, so they take $15$ hours alone.
\dapso{$15$ hours.}}
\gap
\begin{probbox}{4}{\relax}
Two people harvest a field together in $2$ hours. The first alone takes $3$ hours. What fraction of the field does the second harvest each hour?
\end{probbox}
\probsol{
Each hour the second harvests $\dfrac12-\dfrac13=\dfrac16$ of the field.
\dapso{$\dfrac16$ of the field.}}
\gap
\begin{probbox}{5}{\relax}
Three friends tidy the library. Alone, Mai would take $4$ hours, Hung $6$ hours and Lan $12$ hours. How long does it take all three together?
\end{probbox}
\probsol{
Each hour the three do $\dfrac14+\dfrac16+\dfrac1{12}=\dfrac{3+2+1}{12}=\dfrac12$ of the job, so they take $2$ hours.
\dapso{$2$ hours.}}
\gap
\begin{probbox}{6}{\relax}
An inlet fills a tank in $4$ hours; an outlet empties a full tank in $6$ hours. Starting empty with both open, how long until it is full?
\end{probbox}
\probsol{
Each hour the tank gains $\dfrac14-\dfrac16=\dfrac1{12}$, so it is full after $12$ hours.
\dapso{$12$ hours.}}
\gap
\begin{probbox}{7}{\relax}
A tank is already $\dfrac13$ full. A pipe fills an empty tank in $6$ hours. How long will it take the pipe to fill the rest?
\end{probbox}
\probsol{
Still to fill: $1-\dfrac13=\dfrac23$.\par Time: $\dfrac23:\dfrac16=4$ hours.
\dapso{$4$ hours.}}
\gap
\begin{probbox}{8}{\relax}
Pipe I fills a tank in $5$ hours, pipe II in $10$ hours. Starting empty, pipe I runs for $2$ hours and is then replaced by pipe II. How long must pipe II run?
\end{probbox}
\probsol{
Pipe I fills $\dfrac15\times2=\dfrac25$ of the tank, leaving $\dfrac35$.\par Pipe II needs $\dfrac35:\dfrac1{10}=6$ hours.
\dapso{$6$ hours.}}
\gap
\begin{probbox}{9}{\relax}
Team A finishes a job alone in $9$ days, team B in $18$ days. Team A works alone for $3$ days, then team B joins. How many more days do they need together?
\end{probbox}
\probsol{
In $3$ days team A does $\dfrac39=\dfrac13$ of the job, leaving $\dfrac23$.\par
Together they do $\dfrac19+\dfrac1{18}=\dfrac16$ a day, so they need $\dfrac23:\dfrac16=4$ days.
\dapso{$4$ days.}}
\gap
\begin{probbox}{10}{\relax}
Mr Ba finishes a job alone in $10$ days, Mr Tu in $15$ days. They work together for $4$ days, then Mr Ba stops. How long does Mr Tu need to finish?
\end{probbox}
\probsol{
Together they do $\dfrac1{10}+\dfrac1{15}=\dfrac16$ a day, so $\dfrac46=\dfrac23$ in $4$ days, leaving $\dfrac13$.\par
Mr Tu needs $\dfrac13:\dfrac1{15}=5$ days.
\dapso{$5$ days.}}
\gap
\begin{probbox}{11}{\relax}
Two people together finish a job in $4$ hours; the first alone would take $6$ hours. They work together for $2$ hours, then the first stops. How much longer must the second work?
\end{probbox}
\probsol{
The second does $\dfrac14-\dfrac16=\dfrac1{12}$ an hour.\par In $2$ hours together they finish $\dfrac24=\dfrac12$ of the job.\par
The second needs $\dfrac12:\dfrac1{12}=6$ more hours.
\dapso{$6$ hours.}}
\gap
\begin{probbox}{12}{\relax}
Two people do a job together and are paid $1\,800\,000$ dong. Alone, the first would take $6$ days and the second $3$ days. If the pay is shared by the work each did, how much does each get?
\end{probbox}
\probsol{
Each day the first does $\dfrac16$ and the second $\dfrac13=\dfrac26$ of the job, so their work is in the ratio $1:2$.\par
First: $1\,800\,000:3=600\,000$ dong; second: $1\,200\,000$ dong.
\dapso{$600\,000$ dong and $1\,200\,000$ dong.}}
\gap
\begin{probbox}{13}{\relax}
One person finishes a job in $4$ hours, another in $12$ hours. How long do they take together?
\end{probbox}
\probsol{
Each hour together: $\dfrac14+\dfrac1{12}=\dfrac4{12}=\dfrac13$ of the job.\par Together they take $3$ hours.
\dapso{$3$ hours.}}
\gap
\begin{probbox}{14}{\relax}
Pipe I fills a tank in $5$ hours, pipe II in $20$ hours. Starting empty with both open, how long until it is full?
\end{probbox}
\probsol{
Each hour: $\dfrac15+\dfrac1{20}=\dfrac5{20}=\dfrac14$ of the tank.\par Full after $4$ hours.
\dapso{$4$ hours.}}
\gap
\begin{probbox}{15}{\relax}
Two people finish a job together in $8$ days; the first alone takes $12$ days. How long does the second take alone?
\end{probbox}
\probsol{
Each day the second does $\dfrac18-\dfrac1{12}=\dfrac3{24}-\dfrac2{24}=\dfrac1{24}$ of the job.\par Alone: $24$ days.
\dapso{$24$ days.}}
\gap
\begin{probbox}{16}{\relax}
Two pipes fill a tank together in $2$ hours; pipe I alone takes $3$ hours. How long does pipe II take alone?
\end{probbox}
\probsol{
Each hour pipe II fills $\dfrac12-\dfrac13=\dfrac16$ of the tank.\par Alone: $6$ hours.
\dapso{$6$ hours.}}
\gap
\begin{probbox}{17}{\relax}
Three workers alone would take $6$, $10$ and $15$ days. How long do all three take together?
\end{probbox}
\probsol{
Each day: $\dfrac16+\dfrac1{10}+\dfrac1{15}=\dfrac{5+3+2}{30}=\dfrac13$ of the job.\par All three: $3$ days.
\dapso{$3$ days.}}
\gap
\begin{probbox}{18}{\relax}
An inlet fills a tank in $6$ hours; an outlet empties a full tank in $9$ hours. Starting empty with both open, how long until it is full?
\end{probbox}
\probsol{
Each hour: $\dfrac16-\dfrac19=\dfrac3{18}-\dfrac2{18}=\dfrac1{18}$ of the tank.\par Full after $18$ hours.
\dapso{$18$ hours.}}
\gap
\begin{probbox}{19}{\relax}
A tank is $\dfrac14$ full. A pipe fills an empty tank in $8$ hours. How long does it need to fill the rest?
\end{probbox}
\probsol{
Still to fill: $\dfrac34$ of the tank.\par Time: $\dfrac34:\dfrac18=6$ hours.
\dapso{$6$ hours.}}
\gap
\begin{probbox}{20}{\relax}
A tank is $\dfrac23$ full. An outlet empties a full tank in $6$ hours. How long does it take to empty this tank?
\end{probbox}
\probsol{
Each hour it empties $\dfrac16$ of the tank.\par Time: $\dfrac23:\dfrac16=4$ hours.
\dapso{$4$ hours.}}
\gap
\begin{probbox}{21}{\relax}
Team A alone takes $12$ days and team B $6$ days. A works alone for $3$ days, then B joins. How many more days do they need?
\end{probbox}
\probsol{
In $3$ days A does $\dfrac3{12}=\dfrac14$; $\dfrac34$ is left.\par Together they do $\dfrac1{12}+\dfrac16=\dfrac14$ a day.\par Time: $\dfrac34:\dfrac14=3$ days.
\dapso{$3$ days.}}
\gap
\begin{probbox}{22}{\relax}
Pipe I fills a tank in $4$ hours, pipe II in $12$. Starting empty, pipe I runs for $1$ hour, then pipe II is opened too. How much longer until it is full?
\end{probbox}
\probsol{
After $1$ hour pipe I has filled $\dfrac14$; $\dfrac34$ is left.\par Both together: $\dfrac14+\dfrac1{12}=\dfrac13$ an hour.\par Time: $\dfrac34:\dfrac13=\dfrac94$ hours $=2$ h $15$ min.
\dapso{$2$ h $15$ min.}}
\gap
\begin{probbox}{23}{\relax}
Two people together do $\dfrac12$ of a job in $3$ days. The first alone would take $10$ days. How long would the second take alone?
\end{probbox}
\probsol{
Together they do $\dfrac12:3=\dfrac16$ a day.\par The second does $\dfrac16-\dfrac1{10}=\dfrac1{15}$ a day.\par Alone: $15$ days.
\dapso{$15$ days.}}
\gap
\begin{probbox}{24}{\relax}
An does $\dfrac13$ of a job in $2$ hours; Binh does $\dfrac12$ of it in $6$ hours. How long do they take together?
\end{probbox}
\probsol{
An would take $6$ hours for the whole job; Binh $12$ hours.\par Together: $\dfrac16+\dfrac1{12}=\dfrac14$ an hour.\par So $4$ hours.
\dapso{$4$ hours.}}
\gap
\begin{probbox}{25}{\relax}
Each hour pipe I fills $\dfrac16$ of a tank and pipe II $\dfrac13$. What fraction of an empty tank do they fill together in $1$ hour $20$ minutes?
\end{probbox}
\probsol{
$1$ h $20$ min $=\dfrac43$ hours.\par Together: $\dfrac16+\dfrac13=\dfrac12$ of the tank an hour.\par In $\dfrac43$ hours: $\dfrac12\times\dfrac43=\dfrac23$.
\dapso{$\dfrac23$ of the tank.}}
\gap
\begin{probbox}{26}{\relax}
Two teams repair a road together in $6$ days; team I alone takes $10$ days. They work together for $2$ days, then team I leaves. How long does team II need to finish?
\end{probbox}
\probsol{
Team II does $\dfrac16-\dfrac1{10}=\dfrac1{15}$ a day.\par After $2$ days together $\dfrac26=\dfrac13$ is done; $\dfrac23$ is left.\par Team II needs $\dfrac23:\dfrac1{15}=10$ days.
\dapso{$10$ days.}}
\gap
\begin{probbox}{27}{\relax}
A pipe fills a tank in $3$ hours, but because of a leak it takes $4$ hours. How long would the leak take to empty a full tank?
\end{probbox}
\probsol{
Each hour the leak loses $\dfrac13-\dfrac14=\dfrac1{12}$ of the tank.\par It empties the tank in $12$ hours.
\dapso{$12$ hours.}}
\gap
\begin{probbox}{28}{\relax}
Two workers share $2\,000\,000$ dong for a job done together. Alone, the first would take $4$ days and the second $6$. Shared by the work done, how much does each get?
\end{probbox}
\probsol{
Each day the first does $\dfrac14=\dfrac3{12}$ and the second $\dfrac16=\dfrac2{12}$; the work is in the ratio $3:2$.\par First: $2\,000\,000:5\times3=1\,200\,000$ dong; second: $800\,000$ dong.
\dapso{$1\,200\,000$ dong and $800\,000$ dong.}}
\gap
\begin{probbox}{29}{\relax}
Pipe I fills a tank in $30$ minutes, pipe II in $45$ minutes. Together, starting empty, how many minutes do they take?
\end{probbox}
\probsol{
Each minute: $\dfrac1{30}+\dfrac1{45}=\dfrac3{90}+\dfrac2{90}=\dfrac1{18}$ of the tank.\par Full after $18$ minutes.
\dapso{$18$ minutes.}}
\gap
\begin{probbox}{30}{\relax}
Two tractors plough a field together in $4$ days; tractor I alone takes $6$ days. After $2$ days together, tractor I breaks down. How long does tractor II need to finish?
\end{probbox}
\probsol{
Tractor II ploughs $\dfrac14-\dfrac16=\dfrac1{12}$ of the field a day.\par After $2$ days together, half is done.\par Tractor II needs $\dfrac12:\dfrac1{12}=6$ days.
\dapso{$6$ days.}}
\dividerbar{Challenge -- Extension}
\begin{probbox}{31}{\stars{1}}
A pipe fills a sound tank in $2$ hours. Because of a leak, it now takes $2$ hours $30$ minutes. How long would the leak take to empty a full tank?
\end{probbox}
\probsol{
$2$ h $30$ min $=\dfrac52$ hours.\par The pipe adds $\dfrac12$ a tank each hour; with the leak the tank gains only $\dfrac25$ an hour.\par
The leak loses $\dfrac12-\dfrac25=\dfrac1{10}$ a tank each hour, so it empties a full tank in $10$ hours.
\dapso{$10$ hours.}}
\gap
\begin{probbox}{32}{\stars{1}}
A tank is empty. Pipe I fills it in $3$ hours and pipe II in $6$ hours; pipe III empties a full tank in $4$ hours. With all three open, how long until it is full?
\end{probbox}
\probsol{
Each hour the tank gains $\dfrac13+\dfrac16-\dfrac14=\dfrac{4+2-3}{12}=\dfrac14$, so it is full after $4$ hours.
\dapso{$4$ hours.}}
\gap
\begin{probbox}{33}{\stars{1}}
Three workers: alone, A finishes the job in $6$ days, B in $12$ days and C in $12$ days. A and B work together for $2$ days, then C joins. How much longer do the three need?
\end{probbox}
\probsol{
A and B do $\dfrac16+\dfrac1{12}=\dfrac14$ a day, so half the job in $2$ days.\par
All three do $\dfrac14+\dfrac1{12}=\dfrac13$ a day, so they need $\dfrac12:\dfrac13=\dfrac32$ days, i.e. a day and a half.
\dapso{$1.5$ days.}}
\gap
\begin{probbox}{34}{\stars{1}}
The master craftsman finishes a job alone in $12$ hours. He works for $3$ hours, then his assistant works for $4$ hours, and together they have done $\dfrac7{12}$ of the job. How long would they take working together from the start?
\end{probbox}
\probsol{
In $3$ hours the master does $\dfrac3{12}$.\par The assistant's $4$ hours give $\dfrac7{12}-\dfrac3{12}=\dfrac4{12}$, so the assistant does $\dfrac1{12}$ an hour.\par
Together: $\dfrac1{12}+\dfrac1{12}=\dfrac16$ an hour, so $6$ hours.
\dapso{$6$ hours.}}
\gap
\begin{probbox}{35}{\stars{2}}
A and B together finish a job in $12$ days; B and C in $15$ days; A and C in $20$ days.
\begin{cols2}
\pa{a} How long do all three take together? & \pa{b} How long does each take alone?
\end{cols2}
\end{probbox}
\probsol{
\pa{a} Add the three pairs: $\dfrac1{12}+\dfrac1{15}+\dfrac1{20}=\dfrac{5+4+3}{60}=\dfrac15$.\par Each person is counted twice, so all three do $\dfrac15:2=\dfrac1{10}$ a day and take $10$ days.\par
\pa{b} A: $\dfrac1{10}-\dfrac1{15}=\dfrac1{30}$ ($30$ days); B: $\dfrac1{10}-\dfrac1{20}=\dfrac1{20}$ ($20$ days); C: $\dfrac1{10}-\dfrac1{12}=\dfrac1{60}$ ($60$ days).
\dapso{\pa{a} $10$ days; \pa{b} A: $30$, B: $20$, C: $60$ days.}}
\gap
\begin{probbox}{36}{\stars{2}}
Two pipes fill an empty tank. Both together for $2$ hours fill $\dfrac12$ of it. Pipe I for $3$ hours and pipe II for $1$ hour fill $\dfrac7{12}$ of it. How long does each pipe take alone?
\end{probbox}
\probsol{
In $1$ hour both fill $\dfrac12:2=\dfrac14$.\par The second case is ``both for $1$ hour'' plus ``pipe I for $2$ more hours''.\par
So pipe I fills $\dfrac7{12}-\dfrac14=\dfrac13$ in $2$ hours, i.e.\par $\dfrac16$ an hour.\par Pipe II: $\dfrac14-\dfrac16=\dfrac1{12}$ an hour.
\dapso{Pipe I: $6$ hours; pipe II: $12$ hours.}}
\gap
\begin{probbox}{37}{\stars{1}}
A tank is half full. Pipe I fills an empty tank in $6$ hours; pipe II empties a full one in $8$ hours. With both open, how long until it is full?
\end{probbox}
\probsol{
Each hour: $\dfrac16-\dfrac18=\dfrac1{24}$ of the tank.\par Time: $\dfrac12:\dfrac1{24}=12$ hours.
\dapso{$12$ hours.}}
\gap
\begin{probbox}{38}{\stars{1}}
A working $3$ days and B $4$ days do $\dfrac5{12}$ of a job; A working $6$ days and B $2$ days do $\dfrac7{12}$. How long would each take alone?
\end{probbox}
\probsol{
Doubling the first case: A for $6$ days and B for $8$ days do $\dfrac{10}{12}$.\par Compared with the second case, $6$ days of B do $\dfrac{10}{12}-\dfrac7{12}=\dfrac14$, so B does $\dfrac1{24}$ a day.\par A's $3$ days give $\dfrac5{12}-\dfrac4{24}=\dfrac14$, so A does $\dfrac1{12}$ a day.\par A: $12$ days; B: $24$ days.
\dapso{A $12$, B $24$ days.}}
\gap
\begin{probbox}{39}{\stars{1}}
Two people finish a job together in $12$ hours. If the first works $4$ hours and the second $6$ hours, they do $\dfrac25$ of it. How long would each take alone?
\end{probbox}
\probsol{
Working together for $4$ hours gives $\dfrac4{12}=\dfrac13$.\par The second's extra $2$ hours give $\dfrac25-\dfrac13=\dfrac1{15}$, i.e.\par $\dfrac1{30}$ an hour.\par The first: $\dfrac1{12}-\dfrac1{30}=\dfrac1{20}$ an hour.\par First: $20$ h; second: $30$ h.
\dapso{$20$ h and $30$ h.}}
\gap
\begin{probbox}{40}{\stars{1}}
An inlet fills a tank in $4$ hours; an outlet empties a full tank in $6$ hours. Starting empty, the inlet runs for $2$ hours, then the outlet is opened too. How long after the start is the tank full?
\end{probbox}
\probsol{
After $2$ hours the tank is $\dfrac24=\dfrac12$ full.\par With both open it gains $\dfrac14-\dfrac16=\dfrac1{12}$ an hour, needing $\dfrac12:\dfrac1{12}=6$ more hours.\par In all: $2+6=8$ hours.
\dapso{$8$ hours.}}
\gap
\begin{probbox}{41}{\stars{1}}
A and B together finish a job in $6$ days; A alone takes $10$ days. C alone takes twice as long as B. How long do all three take together?
\end{probbox}
\probsol{
B does $\dfrac16-\dfrac1{10}=\dfrac1{15}$ a day, so B takes $15$ days and C $30$ days.\par All three: $\dfrac1{10}+\dfrac1{15}+\dfrac1{30}=\dfrac6{30}=\dfrac15$ a day.\par Together: $5$ days.
\dapso{$5$ days.}}
\gap
\begin{probbox}{42}{\stars{1}}
Three people finish a job together and are paid $3\,600\,000$ dong. Alone, A would take $6$ days, B $12$ and C $4$. Shared by the work done, how much does each get?
\end{probbox}
\probsol{
Each day A does $\dfrac2{12}$, B $\dfrac1{12}$ and C $\dfrac3{12}$; the ratio is $2:1:3$.\par One part: $3\,600\,000:6=600\,000$ dong.\par A: $1\,200\,000$; B: $600\,000$; C: $1\,800\,000$ dong.
\dapso{A $1\,200\,000$; B $600\,000$; C $1\,800\,000$ dong.}}
\gap
\begin{probbox}{43}{\stars{1}}
Running pipe I for $2$ hours and then pipe II for $3$ hours fills a tank. Both together fill it in $2$ hours $24$ minutes. How long does each take alone?
\end{probbox}
\probsol{
$2$ h $24$ min $=\dfrac{12}5$ hours, so together they fill $\dfrac5{12}$ an hour.\par Both for $2$ hours give $\dfrac{10}{12}$; pipe II's extra hour gives the rest, $1-\dfrac{10}{12}=\dfrac16$.\par Pipe II: $6$ h; pipe I fills $\dfrac5{12}-\dfrac16=\dfrac14$ an hour: $4$ h.
\dapso{Pipe I $4$ h, pipe II $6$ h.}}
\gap
\begin{probbox}{44}{\stars{1}}
A alone takes $10$ days and B alone $15$. They start together, but A takes the first $5$ days off while B works throughout. How many days does the job take?
\end{probbox}
\probsol{
In the first $5$ days B does $\dfrac5{15}=\dfrac13$; $\dfrac23$ is left.\par Then both do $\dfrac1{10}+\dfrac1{15}=\dfrac16$ a day, needing $\dfrac23:\dfrac16=4$ days.\par In all: $5+4=9$ days.
\dapso{$9$ days.}}
\gap
\begin{probbox}{45}{\stars{1}}
A tank has two inlets and one outlet. Inlet I fills it in $3$ hours, inlet II in $4$; the outlet empties it in $6$. Starting empty with all three open, how long until it is full?
\end{probbox}
\probsol{
Each hour: $\dfrac13+\dfrac14-\dfrac16=\dfrac{4+3-2}{12}=\dfrac5{12}$ of the tank.\par Time: $1:\dfrac5{12}=\dfrac{12}5$ hours $=2$ h $24$ min.
\dapso{$2$ h $24$ min.}}
\gap
\begin{probbox}{46}{\stars{1}}
Class 5A cleans the schoolyard in $3$ hours, class 5B in $6$. 5A works for $1$ hour, then 5B joins. How much longer do they need?
\end{probbox}
\probsol{
After $1$ hour 5A has done $\dfrac13$; $\dfrac23$ is left.\par Together: $\dfrac13+\dfrac16=\dfrac12$ an hour.\par Time: $\dfrac23:\dfrac12=\dfrac43$ hours $=1$ h $20$ min.
\dapso{$1$ h $20$ min.}}
\gap
\begin{probbox}{47}{\stars{1}}
One pump empties a pond in $10$ hours. After $4$ hours a second pump joins, and $2$ hours later the pond is empty. How long would the second pump take alone?
\end{probbox}
\probsol{
In $4$ hours pump I removes $\dfrac4{10}=\dfrac25$; $\dfrac35$ is left.\par The two remove $\dfrac35$ in $2$ hours: $\dfrac3{10}$ an hour.\par Pump II: $\dfrac3{10}-\dfrac1{10}=\dfrac15$ an hour, so $5$ hours alone.
\dapso{$5$ hours.}}
\gap
\begin{probbox}{48}{\stars{2}}
Pipes I and II together fill a tank in $2$ hours; pipes II and III in $3$ hours; pipe I alone in $3$ hours. How long do all three take together?
\end{probbox}
\probsol{
Pipe II: $\dfrac12-\dfrac13=\dfrac16$ an hour.\par Pipe III: $\dfrac13-\dfrac16=\dfrac16$ an hour.\par All three: $\dfrac13+\dfrac16+\dfrac16=\dfrac23$ an hour.\par Time: $1:\dfrac23=\dfrac32$ hours $=1$ h $30$ min.
\dapso{$1$ h $30$ min.}}
\gap
\begin{probbox}{49}{\stars{2}}
Worker A alone is $6$ hours quicker than worker B alone. Together they finish in $4$ hours. How long does each take alone? (Hint: try values.)
\end{probbox}
\probsol{
Each alone must take more than $4$ hours.\par Try A $5$ h, B $11$ h: $\dfrac15+\dfrac1{11}=\dfrac{16}{55}\ne\dfrac14$.\par Try A $6$ h, B $12$ h: $\dfrac16+\dfrac1{12}=\dfrac3{12}=\dfrac14$ (correct).\par A: $6$ h; B: $12$ h.
\dapso{A $6$ h, B $12$ h.}}
\gap
\begin{probbox}{50}{\stars{2}}
A master craftsman working $8$ hours a day finishes a job in $12$ days. If an assistant who works at half his rate joins him ($8$ hours a day), how many days will the job take?
\end{probbox}
\probsol{
Each day the master does $\dfrac1{12}$ and the assistant $\dfrac1{24}$ of the job.\par Together: $\dfrac1{12}+\dfrac1{24}=\dfrac18$ a day.\par So $8$ days.
\dapso{$8$ days.}}
\clearpage
\titlehead
